\documentclass[11pt]{article}
\usepackage[utf8]{inputenc}
\usepackage[T1]{fontenc}
\usepackage{lmodern}
\usepackage{amsmath,amssymb,amsthm}
\usepackage{geometry}
\usepackage{hyperref}
\usepackage{booktabs}
\geometry{margin=1in}

\title{Turing Instrument Calibration and Validation Program\\
Version 1.0: Experimental Protocol}
\author{M. Drake Stapleton\\AIEN Project}
\date{September 29, 2026}

\begin{document}

\maketitle

\noindent\textit{Supplementary protocol to the preprint ``Computing Machinery and Understanding: From the Imitation Game to the Turing'' (v1.7). The preprint defines the Turing unit and its falsification criteria; this document specifies the ordered experimental program that determines whether the unit is a scientifically useful measure. This protocol is versioned independently of the paper.}

\section{Purpose}

This program is designed to determine whether the \textbf{Turing} is a scientifically useful, reproducible measure of generalized explanatory compression rather than merely a mathematically defined scoring function.

The experimental program must answer, in order:

\begin{enumerate}
\item Does theoretical predictive description length correspond to actual lossless compression?
\item Does the metric reward genuine stochastic structure while rejecting memorization and overfitting?
\item Does additional model structure receive credit only when it earns its complexity?
\item Can the metric distinguish qualitatively different generating mechanisms?
\item Can it reward useful unification across multiple phenomena?
\item Does an explanation remain useful under interventions and distributional changes?
\item Can explanatory models help select better experiments?
\item Are results reproducible under implementation replay and fresh stochastic worlds?
\item Can the associated physical energy accounting be independently closed?
\item Can equivalent explanations be distinguished by explanatory efficiency in Turings per joule?
\item Are scientific conclusions stable under reasonable alternative measurement profiles?
\item Can an independent laboratory reproduce the conclusions?
\end{enumerate}

The program is deliberately structured so that each major hypothesis can fail. Large positive Turing values are not sufficient evidence. The important question is whether the instrument makes correct preregistered discriminations in cases where it is possible for the metric to be wrong.

\section{Core Experimental Principle}

The experiments must be executed as sequential qualification gates. The canonical sequence is:

\begin{verbatim}
CAL-0
Freeze measurement apparatus
->
EXP-001
Connect theoretical description length
to actual lossless compression
->
EXP-002A
Reject false structure in pure randomness
->
EXP-002B
Test when additional structure earns its cost
->
EXP-002C
Distinguish qualitatively different mechanisms
->
EXP-002D
Test explanatory unification
->
Interventional validation
Test frozen laws outside the observational regime
->
EXP-003
Select informative experiments
->
H3
Reproducibility
->
H4
Physical energy accounting
->
H5
Turing yield
->
Sensitivity analysis
->
Independent laboratory replication
\end{verbatim}

A downstream success does not repair an upstream failure. If EXP-001 fails, later experiments may continue for debugging purposes, but they cannot support claims that depend upon a validated Turing measurement.

\section{Measurement Quantity}

For candidate model or program \(M\), baseline \(B\), and sealed held-out data \(D\), calculate:
\[
T(M;B,D)
=
[L(B)+L(D|B)]
-
[L(M)+L(D|M)].
\]
The measurement must account for both \(L(M)\), the description cost of the explanation itself, and \(L(D|M)\), the description length of observations conditional on that explanation.

Predictive codelength is initially calculated as:
\[
L_{\text{ideal}}(D|M)
=
-\sum_t
\log_2
p_M(x_t|x_{<t}).
\]
EXP-001 then determines whether this idealized quantity corresponds closely enough to an actual losslessly encoded artifact to support physical compression claims.

\section{CAL-0: Freeze the Measurement Apparatus}

\subsection{Objective}

Before evaluating any scientific claim, create and freeze a canonical measurement profile:
\begin{verbatim}
Turing-profile-v1.0
\end{verbatim}
The profile may not be modified during the first validation series. A change to the profile creates a new profile version. Results under different profiles must not be silently treated as numerically interchangeable.

\subsection{Required Freeze Order}

The required sequence is:
\begin{verbatim}
experiment designed
->
candidate families specified
->
statistical simulations performed
->
sample size selected
->
pass/fail rules preregistered
->
measurement profile frozen
->
candidate development completed
->
candidate artifact frozen and hashed
->
sealed evaluation data generated/released
->
evaluation performed
->
results retained whether successful or unsuccessful
\end{verbatim}
Candidate development must end before the evaluator obtains sealed test observations. A cryptographic commitment alone does not establish blinding if the candidate has already had access to the hidden data.

\section{\texttt{Turing-profile-v1.0}}

The frozen profile must contain at minimum the following fields.

\subsection{Identity}
\begin{verbatim}
profile_name
profile_version
profile_digest
creation_timestamp
protocol_version
\end{verbatim}

\subsection{Experimental Definition}
\begin{verbatim}
task_definition
hypothesis_definition
primary_endpoint
secondary_endpoints
pass_fail_rule
stopping_rule
\end{verbatim}

\subsection{Baseline}
\begin{verbatim}
baseline_family
baseline_encoding
baseline_parameterization
baseline_description_cost
\end{verbatim}

\subsection{Candidate Encoding}
\begin{verbatim}
candidate_representation
candidate_serialization
shared_background
model_code_method
parameter_precision
dependency_cost
metadata_cost
lookup_table_cost
external_artifact_cost
\end{verbatim}
Particular attention must be paid to \(L(M)\). The experiment must define exactly what information is considered free shared background and what information must be charged to the candidate. No model may hide predictive information in metadata, filenames, headers, preloaded tables, dependencies, or uncharged support files.

\subsection{Data}
\begin{verbatim}
data_generator
development_split
calibration_split
sealed_test_split
replication_split
trajectory_count
trajectory_length
sampling_interval
seed_commitment
\end{verbatim}

\subsection{Blinding}
\begin{verbatim}
holdout_commitment
candidate_environment
evaluator_environment
access_control
test_release_process
candidate_freeze_process
\end{verbatim}

\subsection{Probabilities}
\begin{verbatim}
probability_representation
quantizer
normalization_rule
zero_probability_policy
numeric_precision
\end{verbatim}

\subsection{Coding}
\begin{verbatim}
ideal_codelength_method
reference_coder_A
reference_coder_B
termination_convention
header_accounting
side_information_accounting
\end{verbatim}

\subsection{Statistics}
\begin{verbatim}
uncertainty_method
confidence_level
block_length
bootstrap_method
correlation_handling
number_of_replicates
\end{verbatim}

\subsection{Physical Measurement}
\begin{verbatim}
energy_instrument
sampling_frequency
idle_baseline_method
attribution_method
sensor_uncertainty
time_alignment_method
\end{verbatim}

\subsection{Reproducibility}
\begin{verbatim}
software_versions
compiler_versions
hardware_identity
runtime_digest
candidate_digest
dataset_digest
coder_digest
scorer_digest
\end{verbatim}

\subsection{Sensitivity}
\begin{verbatim}
alternative_baselines
alternative_model_codes
alternative_quantization
alternative_reference_coder
alternative_uncertainty_model
\end{verbatim}

\section{Sample-Size Determination}

The manuscript does not prescribe a fixed sample size. Therefore the experiment must not choose one arbitrarily. Before preregistration:

\begin{enumerate}
\item simulate the proposed hidden processes;
\item simulate the expected range of candidate effect sizes;
\item apply the proposed uncertainty procedure;
\item determine how many independent worlds or trajectories are required to distinguish effects considered scientifically meaningful;
\item freeze that sample size before final sealed-data generation.
\end{enumerate}

The simulation itself must be archived. The sample-size calculation must not be rerun with revised assumptions after sealed outcomes are known unless the original experiment is formally declared inconclusive and a new experiment is preregistered.

\section{Experimental Roles}

At minimum, separate the following roles.

\subsection{Candidate Team}

May access:
\begin{verbatim}
public protocol
development observations
development simulators where permitted
public candidate APIs
\end{verbatim}
Must not access:
\begin{verbatim}
sealed seeds
sealed trajectories
hidden generator parameters
evaluator secrets
future intervention results
\end{verbatim}

\subsection{Evaluator}

Controls:
\begin{verbatim}
sealed observations
sealed parameters
reference coding
scoring
experiment execution
\end{verbatim}

\subsection{Verification Team}

Must independently reconstruct:
\begin{verbatim}
model cost
ideal codelength
actual coded length
Turing value
uncertainty
\end{verbatim}
without depending on the production scoring path.

\subsection{Replication Team}

Uses independent seeds and preferably independent implementations.

\subsection{External Replication Team}

Ultimately generates its own hidden data and implements the measurement protocol independently.

\section{EXP-001: Compression Bridge}

\subsection{Scientific Question}

Does the theoretical predictive codelength used by the Turing correspond to actual physical lossless compression? The experiment must close:
\begin{verbatim}
probability prediction
->
-log2 probability
->
entropy coder
->
real encoded bitstream
->
real compressed size
\end{verbatim}
The current theoretical score alone is insufficient to establish this connection.

\subsection{Candidate Predictors}

Freeze at least:

\begin{center}
\begin{tabular}{lll}
\toprule
ID & Candidate & Purpose \\
\midrule
B0 & uniform & minimal reference \\
B1 & empirical frequency & simple statistical baseline \\
B2 & small Markov & short-range structural baseline \\
B3 & fixed heuristic & domain-specific baseline \\
M1 & explanatory candidate & primary model \\
Mmem & memorizer & adversarial complexity control \\
\bottomrule
\end{tabular}
\end{center}

\section{Probability Artifact}

Every predictor must first produce an immutable probability stream. Each prediction record should include:
\begin{verbatim}
observation_index
context_digest
alphabet
raw_probability_vector
quantized_probability_vector
normalization_receipt
\end{verbatim}
Freeze and hash the probability stream. Both coders must consume exactly this stream. They must not independently invoke the model. This separates model error from coder error.

\section{Dual Reference Coders}

Use two independently implemented coder families. Recommended:
\begin{verbatim}
Coder A:
arithmetic or range coder

Coder B:
ANS or rANS
\end{verbatim}
For each candidate require:
\[
decode(encode(D)) = D.
\]
The decoded sequence must exactly equal the sealed sequence. The two coders need not produce identical binary streams. Their conclusions should agree within the preregistered coding-overhead envelope.

\section{EXP-001 Measurements}

For every candidate compute \(L(M_i)\), \(L_{\text{ideal}}(D|M_i)\), \(L_{\text{coded-A}}(D|M_i)\), and \(L_{\text{coded-B}}(D|M_i)\). Then calculate:
\begin{verbatim}
T_ideal
T_coder_A
T_coder_B
\end{verbatim}
and:
\begin{verbatim}
coder_A_overhead
coder_B_overhead
\end{verbatim}

\section{EXP-001 Adversarial Controls}

Deliberately test:
\begin{verbatim}
corrupted bitstream
truncated bitstream
incorrect profile digest
incorrect model digest
incorrect test-data digest
invalid probability normalization
zero-probability actual event
hidden metadata
candidate/test overlap
unaccounted headers
unaccounted side information
memorized data disguised as program state
\end{verbatim}
Every invalid condition must fail closed.

\section{EXP-001 Acceptance Criteria}

Require:
\begin{verbatim}
both decoders reproduce the sealed data exactly

candidate hashes predate sealed-data release

both coders consume the same probability artifact

actual coded lengths agree with ideal lengths
within preregistered coder overhead

scientifically relevant model ordering remains stable

positive candidate T survives real coding
if preregistered to do so

memorizer does not receive unjustified positive generalized T

independent scorer reproduces the published result

new-seed replication preserves the qualitative conclusion
\end{verbatim}
Gate:
\begin{verbatim}
EXP_001_COMPRESSION_BRIDGE = PASS
\end{verbatim}
or:
\begin{verbatim}
EXP_001_COMPRESSION_BRIDGE = FAIL
\end{verbatim}
A failed result must be retained.

\section{EXP-002A: Randomness Without Invented Structure}

\subsection{Hidden Process}

Generate:
\[
dX_t=\sqrt{2D}\,dW_t.
\]
The evaluator privately samples:
\begin{verbatim}
D
initial condition
trajectory
seed
\end{verbatim}
The scientific question is whether the metric discovers centered stochastic diffusion without inventing deterministic patterns from random fluctuations.

\section{Candidate Families}

Compare:
\begin{verbatim}
weak baseline

simple diffusion

diffusion + drift

mean-reverting model

polynomial/curve fit

high-capacity flexible predictor

memorizer
\end{verbatim}
The correct structural family is deliberately simple. The experiment must provide attractive opportunities for overfitting.

\section{EXP-002A False-Structure Challenge}

Include development trajectories with accidental apparent:
\begin{verbatim}
trend
curvature
local periodicity
reversal
\end{verbatim}
generated purely by noise. Do not remove these runs. They provide the most informative test. A polynomial or flexible predictor may achieve excellent development fit. Its held-out performance must determine whether that structure actually earns explanatory credit.

\section{EXP-002A Measurements}

For each candidate record:
\begin{verbatim}
L(M)
development predictive bits
sealed predictive bits
actual coded bits
T
calibration error
prediction interval coverage
development-to-test degradation
\end{verbatim}
Repeat across hidden variation in:
\begin{verbatim}
D
x0
trajectory length
sampling interval
seed
\end{verbatim}

\section{EXP-002A Preregistered Prediction}

Expected behavior:
\begin{verbatim}
diffusion:
generalizes

diffusion + drift:
does not earn unnecessary drift parameter

mean reversion:
not systematically rewarded

curve fit:
may fit development trajectory
but fails on fresh trajectories

high-capacity predictor:
wins only if predictive savings exceed complexity

memorizer:
fails generalized gain
\end{verbatim}
Gate:
\begin{verbatim}
EXP_002A_RANDOMNESS_DISCRIMINATION = PASS
\end{verbatim}

\section{EXP-002B: Occam Transition}

\subsection{Hidden Process}

Generate:
\[
dX_t=\mu\,dt+\sqrt{2D}\,dW_t.
\]
The question is not merely whether drift can be detected. The question is when its predictive value becomes large enough to pay for the additional parameter.

\section{Drift Sweep}

Preregister a sweep:
\begin{verbatim}
mu = 0
very small
small
medium
large
\end{verbatim}
Actual numeric values must be chosen through power simulation and frozen before final data generation. Where appropriate vary:
\begin{verbatim}
sign(mu)
D
x0
trajectory duration
\end{verbatim}

\section{Primary Comparison}

Compare:
\begin{verbatim}
M0:
diffusion

M1:
diffusion + drift
\end{verbatim}
Calculate \(\Delta L(M)\), \(\Delta L(D|M)\), and \(\Delta T\). At sufficiently small drift:
\[
\Delta L(D|M)<\Delta L(M),
\]
so drift should not earn its cost. At sufficiently strong drift:
\[
\Delta L(D|M)>\Delta L(M),
\]
so the richer explanation should win.

\section{Occam Curve}

Report:
\begin{verbatim}
effect size
delta model bits
delta predictive bits
delta T
uncertainty interval
probability richer model wins
\end{verbatim}
Estimate the region where the richer model crosses from unjustified to justified. The \(\mu=0\) condition is mandatory. Systematic preference for drift at \(\mu=0\) is a failure requiring investigation.

Gate:
\begin{verbatim}
EXP_002B_OCCAM_TRANSITION = PASS
\end{verbatim}

\section{EXP-002C: Mechanism Discrimination}

\subsection{Hidden Process}

Generate:
\[
dX_t=-\theta X_t\,dt+\sigma\,dW_t.
\]
This experiment tests a qualitatively different mechanism: state-dependent restoring dynamics.

\section{Candidate Families}

Compare:
\begin{verbatim}
Brownian motion
Brownian motion + constant drift
Ornstein-Uhlenbeck / mean reversion
flexible black-box predictor
trajectory-fitting model
\end{verbatim}

\section{Observational Phase}

Develop candidates on a standard observational distribution. Freeze all candidate artifacts before intervention conditions are revealed.

\section{Intervention Phase}

Evaluate from:
\begin{verbatim}
near-equilibrium initial state

moderately displaced initial state

far-displaced initial state
\end{verbatim}
A genuine mean-reverting explanation predicts that expected restoring behavior changes systematically with state. A trajectory-fitting predictor may fail.

\section{Explanatory Evidence Rule}

Do not call successful passive compression explanatory by itself. The stronger interpretation requires evidence such as:
\begin{verbatim}
invariance

interventional prediction

counterfactual success

mechanistic structure
\end{verbatim}
A particularly strong result is:
\begin{verbatim}
model frozen
->
initial condition changed
->
new regime unseen during fitting
->
same compact law predicts outcomes
->
compression advantage survives
\end{verbatim}
Gate:
\begin{verbatim}
EXP_002C_MECHANISM_DISCRIMINATION = PASS
\end{verbatim}

\section{EXP-002D: Unification}

Construct a family containing:
\begin{verbatim}
pure diffusion
drifted diffusion
mean reversion
different D
different mu
different theta
different time steps
different initial states
\end{verbatim}
Compare two architectures.

\subsection*{Specialists}

Separate predictor for each environment. Charge every model independently.

\subsection*{Unified Theory}

One common framework plus small environment-specific parameters. Charge:
\begin{verbatim}
shared theory
parameterization
environment parameters
selection metadata
\end{verbatim}
The source explicitly proposes testing whether an explanation can compress multiple hidden worlds and earn aggregate gain through genuine reuse.

\section{Unification Scaling Curve}

Evaluate:
\begin{verbatim}
1 environment
2 environments
4 environments
8 environments
...
\end{verbatim}
Observe whether the initial cost of a common abstraction eventually becomes cheaper than maintaining separate specialist theories. Also include a heterogeneous negative-control collection where forced unification should not win. Success requires both:
\begin{verbatim}
reward useful unification

reject artificial unification
\end{verbatim}
Gate:
\begin{verbatim}
EXP_002D_UNIFICATION = PASS
\end{verbatim}

\section{EXP-003: Active Scientific Agency}

This experiment asks whether learned explanations help determine what evidence should be collected next. It is potentially the most important step because it tests whether explanatory representations improve action rather than merely retrospective scoring.

\section{Ambiguous Starting State}

Construct observations compatible with at least two hypotheses. Example:
\begin{verbatim}
H0:
free diffusion

H1:
mean reversion
\end{verbatim}
The initial evidence must not trivially distinguish them.

\section{Permitted Experiments}

Define a finite preregistered action space such as:
\begin{verbatim}
change initial displacement
increase observation duration
change sample frequency
request another replicate
\end{verbatim}
Every permitted intervention receives an explicit cost.

\section{Candidate Decision Loop}

At every round:
\begin{verbatim}
current observations
->
competing hypotheses
->
predictive distribution under each possible intervention
->
expected information gain
->
rank interventions
->
commit chosen intervention
->
evaluator executes experiment
->
result revealed
->
hypotheses rescored
\end{verbatim}
All predictions and choices must be committed before the result is revealed.

\section{Passive Baselines}

Freeze:
\begin{verbatim}
fixed observation schedule
random legal intervention
regular non-adaptive schedule
\end{verbatim}
Use identical discrimination thresholds.

\section{Primary EXP-003 Endpoint}

Measure \(N_{\text{active}}\) against \(N_{\text{passive}}\), where \(N\) is the number of observations required to cross a preregistered discrimination threshold. Also record:
\begin{verbatim}
time-to-discrimination
experimental operations
T accumulated
energy used
\end{verbatim}
Before H4 passes, energy results must be labeled provisional.

Gate:
\begin{verbatim}
EXP_003_ACTIVE_EXPERIMENT_SELECTION = PASS
\end{verbatim}
only if the active strategy uses meaningfully fewer observations or lower preregistered experimental cost than passive baselines.

\section{Counterfactual Experiment Audit}

After each completed active run, use the hidden simulator offline to calculate what would have happened under the other permitted interventions. The live candidate never receives these counterfactual results. This allows evaluation of whether its selected experiment was actually informative relative to alternatives.

\section{H3: Reproducibility}

Reproducibility contains two separate requirements.

\subsection{H3-A: Replay}

Using identical:
\begin{verbatim}
candidate
profile
data
seed
\end{verbatim}
an independent implementation must reconstruct the same result within numerical tolerance.

\subsection{H3-B: Statistical Robustness}

Using:
\begin{verbatim}
new seeds
new trajectories
\end{verbatim}
the exact Turing value may change. The preregistered scientific conclusion should remain.

Gate:
\begin{verbatim}
H3_REPLAY_REPRODUCIBILITY = PASS

H3_STATISTICAL_ROBUSTNESS = PASS
\end{verbatim}

\section{H4: Energy Attribution}

T/J is not validated until the physical energy-accounting system itself is validated. Measure \(E_0\) for idle/background operation, \(E_A\) for workload A, \(E_B\) for workload B, and \(E_{AB}\) for concurrent execution.

\section{Interaction Energy}

Calculate:
\[
I_{AB}
=
E_{AB}-E_A-E_B+E_0.
\]
Do not assume this value is zero. Concurrent workloads may exhibit:
\begin{verbatim}
contention
shared efficiency
thermal interaction
memory interaction
scheduler interaction
\end{verbatim}

\section{Energy Categories}

Allow accounting to distinguish:
\begin{verbatim}
A-attributed
B-attributed
shared
interaction
unattributed
measurement uncertainty
\end{verbatim}
Do not force every joule into one workload merely to make the accounting look clean.

\section{H4 Test Matrix}

Measure:
\begin{verbatim}
A alone
B alone
A+B

CPU-heavy + CPU-heavy

GPU-heavy + GPU-heavy

CPU-heavy + GPU-heavy

memory-bandwidth contention

light/heavy combinations

idle background
\end{verbatim}
Randomize run order where practical. Repeat across representative thermal and load states.

\section{Conservation Requirement}

For each test require:
\begin{verbatim}
measured energy
~=
attributed
+ shared
+ interaction
+ unattributed
\end{verbatim}
within preregistered sensor and integration uncertainty.

Gate:
\begin{verbatim}
H4_ENERGY_ACCOUNTING_CLOSURE = PASS
\end{verbatim}
Until this passes, validated T/J claims are prohibited.

\section{H5: Turing Yield}

Once both T and energy are independently validated, compare semantically equivalent realizations. Example:
\begin{verbatim}
reference CPU
optimized CPU
accelerated realization
\end{verbatim}
The candidates must produce equivalent predictive distributions inside a preregistered tolerance. Require:
\[
T_A \approx T_B.
\]
Then measure:
\[
Y_A = \frac{T_A}{E_A},
\qquad
Y_B = \frac{T_B}{E_B}.
\]

\section{Separate Evaluation and Discovery Yield}

Report:

\subsection*{Evaluation Yield}
\[
T/J_{\text{evaluation}}
\]
Energy required to execute an already-discovered explanation.

\subsection*{Discovery Yield}
\[
T/J_{\text{discovery}}
\]
Energy expended discovering that explanation. A system that requires enormous discovery energy but cheap inference must not be advertised solely by inference efficiency.

Gate:
\begin{verbatim}
H5_TURING_YIELD = PASS
\end{verbatim}
requires at least one semantically equivalent pair whose T/J difference exceeds the combined uncertainty envelope.

\section{Sensitivity Analysis}

Every major successful result must be deliberately attacked. Vary reasonable preregistered alternatives for:
\begin{verbatim}
baseline
model-code scheme
parameter precision
probability quantization
reference coder
trajectory length
holdout regime
seed
uncertainty method
block-resampling assumptions
\end{verbatim}
The numerical Turing value may change. The core question is whether scientific conclusions change.

\section{Rank-Reversal Analysis}

For every meaningful reversal classify the cause as:
\begin{verbatim}
decision-boundary effect
sample-size insufficiency
baseline sensitivity
model-code sensitivity
quantization sensitivity
statistical instability
implementation defect
fundamental metric sensitivity
\end{verbatim}
Do not average reversals away.

Gate:
\begin{verbatim}
TURING_PROFILE_SENSITIVITY = PASS
\end{verbatim}
only under a preregistered qualitative-stability criterion.

\section{Independent Laboratory Replication}

The framework should not be considered strongly validated until at least one external laboratory independently reproduces a subset of the program. The external laboratory should:
\begin{verbatim}
generate its own hidden observations

control its own blinding

implement the published measurement profile

write its own scorer

preferably write its own entropy coder

run EXP-001

run at least one full EXP-002 family

test one reasonable alternative profile
\end{verbatim}
The central replication question is not:
\begin{verbatim}
Did both laboratories report exactly the same T?
\end{verbatim}
It is:
\begin{verbatim}
Did the same scientific candidates
win and lose for the same reasons?
\end{verbatim}

\section{Experimental Artifact Standard}

Each experiment should produce an immutable bundle. Recommended structure:
\begin{verbatim}
experiment/
|-- preregistration.json
|-- profile.json
|-- profile.sha256
|-- environment.json
|-- candidate_manifest.json
|-- candidates/
|-- development_data/
|-- sealed_data_commitment/
|-- probability_streams/
|-- coded_artifacts/
|-- raw_measurements/
|-- energy/
|-- statistical_analysis/
|-- independent_verification/
|-- sensitivity/
|-- replication/
|-- final_receipt.json
|-- REPORT.md
\end{verbatim}
The sealed observations themselves may remain private until the declared release condition.

\section{Canonical Experiment Receipt}

Every experiment should emit a receipt containing:
\begin{verbatim}
experiment_id
experiment_version

profile_digest

candidate_digests
baseline_digest

development_dataset_digest
sealed_dataset_digest

coder_digests
scorer_digests

software_digest
hardware_identity

start_time
end_time

primary_endpoint
result

T
uncertainty

energy
energy_uncertainty

pass_fail

preregistered_gate_digest

artifact_root

verifier_signature
\end{verbatim}

\section{Statistical Analysis Principles}

\subsection*{Independence}
Do not treat correlated observations within a trajectory as independent samples. Use trajectory-level replication or block resampling as defined by the profile.

\subsection*{Multiple Comparisons}
If many candidate families or parameters are evaluated, define the correction or hierarchical analysis before evaluation.

\subsection*{Failed Runs}
A failed computation may be excluded only under preregistered technical-failure criteria. Scientifically inconvenient outcomes may not be discarded.

\subsection*{Stopping}
No optional stopping unless explicitly modeled and preregistered.

\subsection*{Seeds}
Seeds must be committed before final evaluation.

\subsection*{Uncertainty}
Report uncertainty around:
\begin{verbatim}
T
delta T
model ordering
crossover point
T/J
energy attribution
active/passive efficiency
\end{verbatim}
rather than publishing only point estimates.

\section{Failure Policy}

A failed gate should produce:
\begin{verbatim}
FAIL receipt
raw evidence
diagnosis
corrective hypothesis
\end{verbatim}
The next version becomes a new experiment. Example:
\begin{verbatim}
EXP-002B-v1
FAIL

EXP-002B-v2
new preregistration
new sealed worlds
\end{verbatim}
Do not silently rerun a failed experiment until success.

\section{Claim Ladder}

Scientific claims should be limited to the highest completed evidence level.

\begin{center}
\begin{tabular}{lll}
\toprule
Level & Required Evidence & Permitted Interpretation \\
\midrule
0 & Definition only & mathematically specified statistic \\
1 & EXP-001 & reproducibly connected to actual lossless compression in tested regime \\
2 & EXP-002A/B/C & discriminates supported from unsupported stochastic structure \\
3 & interventions + EXP-002D & evidence for invariant/unifying explanatory structure \\
4 & EXP-003 & explanatory models improve experimental choice \\
5 & H4 & physical energy attribution validated \\
6 & H5 & explanatory efficiency measurable as T/J in tested regime \\
7 & sensitivity & conclusion stable across reasonable profile choices \\
8 & independent replication & independently reproducible metrology \\
9 & multiple unrelated domains & evidence for broader applicability \\
\bottomrule
\end{tabular}
\end{center}

\section{Final Qualification Matrix}

At the conclusion of the initial calibration program publish:
\begin{verbatim}
TURING_PROFILE_V1_FROZEN
PASS / FAIL

EXP_001_COMPRESSION_BRIDGE
PASS / FAIL

EXP_002A_RANDOMNESS_DISCRIMINATION
PASS / FAIL

EXP_002B_OCCAM_TRANSITION
PASS / FAIL

EXP_002C_MECHANISM_DISCRIMINATION
PASS / FAIL

EXP_002D_UNIFICATION
PASS / FAIL

EXP_003_ACTIVE_EXPERIMENT_SELECTION
PASS / FAIL

H3_REPLAY_REPRODUCIBILITY
PASS / FAIL

H3_STATISTICAL_ROBUSTNESS
PASS / FAIL

H4_ENERGY_ACCOUNTING_CLOSURE
PASS / FAIL

H5_TURING_YIELD
PASS / FAIL

TURING_PROFILE_SENSITIVITY
PASS / FAIL

INDEPENDENT_REPLICATION
PASS / FAIL
\end{verbatim}

\section{Recommended Execution Order}

The engineering program should proceed in the following order:
\begin{verbatim}
1. Freeze model encoding and L(M)

2. Build probability-stream artifact format

3. Build independent arithmetic/range coder

4. Build independent ANS/rANS coder

5. Build independent Turing verifier

6. Complete CAL-0

7. Run EXP-001

8. Repeat EXP-001 on independent seeds

9. Run EXP-002A

10. Run EXP-002B

11. Run EXP-002C

12. Run interventional EXP-002C

13. Run EXP-002D

14. Run H3 internal replication

15. Run EXP-003

16. Validate physical energy measurement with H4

17. Run H5 Turing-yield experiments

18. Execute profile-sensitivity matrix

19. Freeze the complete calibration record

20. Hand the protocol to an independent laboratory
\end{verbatim}

\section{Immediate Milestone}

The first milestone is deliberately narrow:

\begin{quote}
Produce one fully reproducible EXP-001 bundle in which frozen predictors emit immutable probability streams, two independently implemented lossless coders produce physical compressed artifacts, a separate verifier reconstructs all Turing values, and the explanatory candidate and memorization control behave according to the preregistered prediction.
\end{quote}

Nothing later in the program is more important until that succeeds.

\section{Final Scientific Standard}

The validation program should not be judged by whether the experiments generate impressive Turing numbers. It should be judged by whether the instrument survives cases designed to falsify it.

A persuasive calibration record would show that the metric:
\begin{verbatim}
rewards correct diffusion

rejects invented drift when mu = 0

rewards real drift only after it earns its description cost

rejects deterministic fits to stochastic fluctuations

recognizes mean reversion only when it exists

continues to predict after intervention

rewards useful theoretical unification

rejects artificial unification

helps choose discriminating experiments

reproduces under new trajectories

reproduces under independent implementation

survives reasonable profile changes

accounts physically for its energy use

distinguishes equivalent explanations by T/J

and reproduces in an independent laboratory
\end{verbatim}

If the program achieves those outcomes under frozen, blinded, preregistered conditions, the Turing would have moved from a defined theoretical instrument to an empirically calibrated measurement framework. If it fails any of them, that failure identifies exactly which scientific interpretation has not yet been earned.

\end{document}
